\section{Graph terms}\label{app:graphs}

The terms of graph theory that the sections above use, each with the place where it is used. Graph,
path, connected component and directed graph are defined in Appendix~\ref{gi-app:graphs} of the note.

\paragraph{Shape} (Section~\ref{sec:chains}). For a connected graph of $n$ nodes and $m$ edges:
\begin{itemize}
  \item \textbf{Density} $m/\binom{n}{2}$: the share of the node pairs joined by an edge.
  \item \textbf{Distance} between two nodes: the fewest edges of a path joining them.
        \textbf{Diameter}: the largest distance. In the chain 1--2--3 the diameter is 2; in a
        triangle, 1.
  \item \textbf{Complete graph}: every pair joined, density 1, diameter 1. \textbf{Tree}: $m=n-1$,
        the fewest edges that keep $n$ nodes connected, with no closed path; a \emph{chain} is a tree
        with diameter $n-1$, a \emph{star} one with diameter 2.
\end{itemize}

\paragraph{Order} (Section~\ref{sec:leaders}). In a directed graph:
\begin{itemize}
  \item a \textbf{cycle} is a path that follows the arrows back to its start (A$\to$B$\to$C$\to$A);
        a graph with none is \textbf{acyclic}, and its nodes can be put in a line with every arrow
        pointing forward;
  \item a \textbf{source} has no arrow coming in: here, a pilot with nobody ahead, a leader;
  \item the \textbf{depth} of an acyclic graph is its longest chain of arrows: 1 for a star, $n-1$
        for a line.
\end{itemize}
A directed triangle is \emph{transitive} (A$\to$B, B$\to$C, A$\to$C) or \emph{cyclic}; with arrows
drawn at random, 2 of its 8 orientations are cyclic.

\paragraph{Configuration model} (Section~\ref{sec:partners}). A random graph with the
\emph{degrees} (the number of edges of each node) of a given one: each node gets as many half-edges
as its degree, and the half-edges are paired at random~\cite{newman2010}. It keeps how many
partners each pilot had and changes only who they were.
